% 01-scope.tex — Clause 1: Scope
\chapter{Scope}
\label{sec:scope}
\index{scope (of document)}

\pnum
This document specifies the Lain programming language.  It defines the
syntax, the semantic rules, and the constraints that every conforming Lain
implementation shall enforce.  It serves as the single authoritative
reference for language users and implementors alike.

\pnum
Lain is a statically-typed, compiled programming language designed for
systems programming, embedded systems, and safety-critical software.  It
achieves memory safety, type safety, and resource safety entirely at compile
time with zero runtime overhead.

\pnum
This specification covers:

\begin{itemize}
  \item lexical structure and grammar (\S\,\ref{sec:05-lexical});
  \item type system and type inference (\S\,\ref{sec:07-types});
  \item variable declarations, scoping, and initialization
        (\S\,\ref{sec:06-basics}, \S\,\ref{sec:10-declarations});
  \item expressions, operators, and evaluation order
        (\S\,\ref{sec:08-expressions});
  \item statements and control flow (\S\,\ref{sec:09-statements});
  \item functions, procedures, and purity (\S\,\ref{sec:12-functions});
  \item ownership, borrowing, and the linear type system
        (\S\,\ref{sec:11-ownership});
  \item value range constraints and static verification
        (\S\,\ref{sec:13-vra});
  \item compile-time evaluation and generics (\S\,\ref{sec:14-generics});
  \item pattern matching and exhaustiveness checking
        (\S\,\ref{sec:15-match});
  \item module system (\S\,\ref{sec:16-modules});
  \item C interoperability (\S\,\ref{sec:17-interop});
  \item unsafe code boundaries (\S\,\ref{sec:18-unsafe});
  \item memory model (\S\,\ref{sec:19-memory});
  \item diagnostic messages (Annex~B);
  \item standard library (\S\,\ref{sec:20-stdlib}).
\end{itemize}

\pnum
This specification does \emph{not} cover:

\begin{itemize}
  \item the internal architecture of any particular compiler implementation;
  \item build system integration or tooling;
  \item IDE or editor support;
  \item operating system interfaces beyond what the standard library provides.
\end{itemize}

\section{Compilation model}
\label{sec:scope.compilation}
\index{compilation model}

\pnum
Lain uses a source-to-C transpilation model.  A Lain source file (extension
\texttt{.ln}) is processed by a conforming implementation to produce C99
source code, which is subsequently compiled by a C99 compiler to produce an
executable.  The observable behavior of the program is defined by this
specification; the form of the intermediate C99 code is
\idbehav{the generated C99 code need not be human-readable}.

\pnum
A conforming implementation shall perform the following translation phases
in the following order:

\begin{enumerate}
  \item \textbf{Lexical analysis} --- source text is converted to a token
        stream as described in \S\,\ref{sec:05-lexical}.
  \item \textbf{Parsing} --- the token stream is reduced to an abstract
        syntax tree according to the grammar in Annex~A.
  \item \textbf{Name resolution} --- identifiers are resolved to
        declarations as described in \S\,\ref{sec:06-basics}.
  \item \textbf{Type checking} --- types are inferred and verified as
        described in \S\,\ref{sec:07-types} through
        \S\,\ref{sec:10-declarations}.
  \item \textbf{Exhaustiveness checking} --- pattern matching completeness
        is verified as described in \S\,\ref{sec:15-match}.
  \item \textbf{Use analysis (NLL)} --- last-use points are computed for
        all variables.
  \item \textbf{Linearity and borrow checking} --- ownership and borrowing
        rules are verified as described in \S\,\ref{sec:11-ownership}.
  \item \textbf{Value range analysis} --- bounds safety and numeric
        constraints are verified as described in \S\,\ref{sec:13-vra}.
  \item \textbf{Code emission} --- well-formed programs are translated to
        C99.
\end{enumerate}

\pnum
A program that completes all phases without a diagnostic message is
\textbf{well-formed} and its translation output is defined.

\begin{note}
  Phases~3--8 collectively form \emph{semantic analysis}.  An
  implementation may interleave these phases as long as the observable
  behavior is equivalent to executing them in the order listed.
\end{note}
